\documentclass[varwidth=15cm, border=16pt]{standalone}

\usepackage[utf8]{inputenc}
\usepackage[T2A]{fontenc}
\usepackage[russian]{babel}
\usepackage{amsmath, amssymb}
\usepackage{xcolor}
\usepackage{tikz}
\usetikzlibrary{arrows.meta, positioning}
\usepackage[most]{tcolorbox}

% ── Palette ─────────────────────────────────────────────────────────────────
\definecolor{cblue}{HTML}{2A529A}
\definecolor{cbluelt}{HTML}{EBF1FB}
\definecolor{camber}{HTML}{B37012}
\definecolor{camberlt}{HTML}{FEF5E5}
\definecolor{cdark}{HTML}{2B2D42}
\definecolor{cgray}{HTML}{7A8A9E}
\definecolor{cline}{HTML}{C8D4E4}

\begin{document}

\begin{tcolorbox}[
  enhanced,
  colback=white,
  colframe=cline,
  arc=9pt,
  boxrule=1.2pt,
  left=14pt, right=14pt, top=14pt, bottom=14pt,
]

%% ── Заголовок ────────────────────────────────────────────────────────────
\noindent{\large\bfseries\color{cblue} Q-функция и TD-ошибка}

\vspace{5pt}
{\color{cline}\rule{\linewidth}{0.7pt}}
\vspace{8pt}

%% ── Блок Q-функции ───────────────────────────────────────────────────────
\begin{tcolorbox}[
  enhanced,
  colback=cbluelt,
  colframe=cblue,
  arc=5pt,
  boxrule=0.8pt,
  coltitle=white,
  colbacktitle=cblue,
  fonttitle=\small\bfseries,
  title=Что такое Q-функция?,
  toptitle=4pt, bottomtitle=4pt,
  top=9pt, bottom=7pt, left=10pt, right=10pt,
]
\[
  Q(s,\,a)
  \;=\;
  \mathbb{E}\!\bigl[G_t \;\big|\; s_t = s,\; a_t = a\bigr],
  \qquad
  G_t \;=\; \sum_{k=0}^{\infty} \gamma^k\, r_{t+k}
\]
\vspace{1pt}
{\small \textbf{Смысл:} $Q(s,a)$ — ожидаемая суммарная дисконтированная
награда, если в состоянии $s$ выбрать действие $a$ и затем следовать
текущей стратегии.}
\end{tcolorbox}

\vspace{8pt}

%% ── Блок TD-ошибки ───────────────────────────────────────────────────────
\begin{tcolorbox}[
  enhanced,
  colback=camberlt,
  colframe=camber,
  arc=5pt,
  boxrule=0.8pt,
  coltitle=white,
  colbacktitle=camber,
  fonttitle=\small\bfseries,
  title=Что такое TD-ошибка?,
  toptitle=4pt, bottomtitle=4pt,
  top=9pt, bottom=7pt, left=10pt, right=10pt,
]
\[
  \delta
  \;=\;
  \underbrace{r + \gamma \max_{a'} Q(s',a')}_{\displaystyle\text{новая оценка}}
  \;-\;
  \underbrace{Q(s,\,a)}_{\displaystyle\text{старая оценка}}
\]
\vspace{1pt}
{\small \textbf{Смысл:} TD-ошибка — «удивление» агента: насколько
реальный ход меняет его прогноз о ценности выбранного действия.}
\end{tcolorbox}

\vspace{12pt}

%% ── Диаграмма перехода ───────────────────────────────────────────────────
\begin{center}
\begin{tikzpicture}[
  font=\small,
  box/.style={
    draw=cblue!65, fill=cbluelt, rounded corners=5pt,
    minimum width=2.45cm, minimum height=0.88cm,
    align=center, line width=0.75pt, text=cdark,
  },
  arr/.style={-{Stealth[length=5pt, width=4pt]}, thick, cgray},
]
  \node[box] (s)   {$s$\\[2pt]{\scriptsize\color{cgray}состояние}};
  \node[box, right=1.05cm of s]  (a)  {$a$\\[2pt]{\scriptsize\color{cgray}действие}};
  \node[box, right=1.05cm of a]  (r)  {$r$\\[2pt]{\scriptsize\color{cgray}награда}};
  \node[box, right=1.05cm of r]  (sp) {$s'$\\[2pt]{\scriptsize\color{cgray}след.\ сост.}};

  \draw[arr] (s)  -- (a);
  \draw[arr] (a)  -- (r);
  \draw[arr] (r)  -- (sp);

  \node[below=7pt of a, align=center, font=\scriptsize] {
    {\color{cblue}\bfseries $Q(s,a)$}\\[1pt]{\color{cgray}старая оценка}
  };
  \node[below=7pt of sp, align=center, font=\scriptsize] {
    {\color{camber}\bfseries $r+\gamma\max_{a'}Q(s',a')$}\\[1pt]{\color{cgray}новая цель}
  };
\end{tikzpicture}
\end{center}

\vspace{9pt}
{\color{cline}\rule{\linewidth}{0.7pt}}
\vspace{5pt}

%% ── Итоговая заметка ─────────────────────────────────────────────────────
{\small\color{cgray}%
\textbf{\color{cblue}Связь с PER:}
чем больше $|\delta|$, тем «удивительнее» переход —
тем чаще он попадает в обучающий батч.}

\end{tcolorbox}

\end{document}
